% Online Resource 1 -- technical notes and Tables S1-S6
% Build: pdflatex ESM_1 (twice). The tables between the markers are written by make_supplement.py.
\documentclass[10pt,a4paper]{article}
\usepackage[T1]{fontenc}
\usepackage{lmodern}
\usepackage[margin=2cm]{geometry}
\usepackage{amsmath,amssymb}
\usepackage{booktabs,longtable,array}
\usepackage[hidelinks]{hyperref}
\renewcommand{\thetable}{S\arabic{table}}
\pagestyle{plain}
\setlength{\LTcapwidth}{\textwidth}

\begin{document}

\begin{center}
{\Large\bfseries Online Resource 1}\\[0.8em]
{\large Birthday Polynomials and the Boundaries of Newton's Method}\\[0.4em]
\emph{The Mathematical Intelligencer}\\[0.8em]
Americo Cunha Jr.\\[0.2em]
{\small Laborat\'orio Nacional de Computa\c{c}\~ao Cient\'ifica, Petr\'opolis, Brazil\\
Universidade do Estado do Rio de Janeiro, Rio de Janeiro, Brazil\\
\texttt{americo@lncc.br}}
\end{center}

\bigskip

This Online Resource contains technical notes on the computations of the article and
the data behind the numbers it quotes. The 38
dates are those of the figures, those mentioned in the text and in the box on conventions,
and eight further dates from the history of science and of Brazil. Dates are written
day/month/year. Table~\ref{tab:s1} gives the polynomial of each date and its
multiplicities. Table~\ref{tab:s2} gives the window of the pictures and the data for the
Girard heuristic. Table~\ref{tab:s3} gives the classification of the free critical points
and the result of the numerical test of the regular ones. Table~\ref{tab:s4} collects the
runs of the Lyapunov indicator for 20/11/1924, and Table~\ref{tab:s5} the verification of
the critical-orbit intervals for that date. Table~\ref{tab:s6} classifies the points of
each fractal picture of the article.

Exact quantities are the polynomials, degrees, multiplicities, and numbers of distinct
roots (Table~\ref{tab:s1}), the classification of the critical points
(Table~\ref{tab:s3}), and the checks of Table~\ref{tab:s5}, whose enclosures are validated
by interval arithmetic. Decimal numbers are rounded approximations, given with three
decimal places (four for the Lyapunov runs); unless stated otherwise, they were obtained
in double-precision floating point. The convergence recorded in Table~\ref{tab:s3} is a
finite numerical criterion, not a proof of convergence. The only nonconvergence in that
table, that of the real critical point of 20/11/1924, is proved in the article
(Proposition~3).

{\sloppy The same data are in the files \texttt{results/dates.csv},
\texttt{results/lyapunov.json}, \texttt{results/orbit\_intervals\_certificate.json},
\texttt{results/trapping\_certificate.json}, and \texttt{results/picture\_classification.json} of the code repository of the article, \url{https://github.com/americocunhajr/BirthdayFractal} (folder \texttt{paper}), which contains the programs that produce them. Its \texttt{README.md}
describes how to reproduce the data and these tables, and the program
\texttt{audit.py} recomputes the numbers quoted in the article and in these
notes and compares them with the published values.\par}

\section*{Technical notes}

These notes give the computational details that the article summarizes.
Section, figure, proposition, and corollary numbers refer to the article.

\emph{Arithmetic and software.} The computations are of three kinds. The
pictures, the test of the critical orbits, and the Lyapunov indicator use
double-precision floating point (or higher precision where stated). The
classification of roots, multiplicities, and critical points, and the
certificate for Proposition~3, use exact rational arithmetic. The multiplier of
the period-two orbit of Corollary~1 and the verification of the critical-orbit
intervals use interval arithmetic with outward rounding. The computations were
done with Python~3.13.16, NumPy~2.5.3, Matplotlib~3.11.2, Pillow~12.3.0,
mpmath~1.3.0, and SymPy~1.14.0; they are deterministic and involve no random
numbers. The interactive website mentioned in Section~1,
\url{https://americocunha.org/BirthdayFractal/}, is a separate, simplified
implementation in JavaScript, with its own colors and resolutions; all figures
and numbers of the article come from the programs of the code repository.

\emph{Pictures.} The parameters that affect the classification are the grid,
$3200\times3200$ points reduced to $1600\times1600$ pixels; the stopping
criterion, a Newton step $\Delta=p(z_k)/p'(z_k)$ with $|\Delta|<10^{-10}R$,
where $R$ is the half-width of the window, or 300 iterations; and the
classification radius $10^{-4}R$. Each point falls into one of four classes:
(a)~\emph{classified}, if it stopped by the step criterion within $10^{-4}R$ of
a root, which defines its color; (b)~\emph{non-finite}, if the step or the
iterate ceased to be a finite floating-point number; (c)~\emph{unclassified
after 300 iterations}; (d)~\emph{stopped far from the roots}, if it stopped by
the step criterion farther than $10^{-4}R$ from every root. Classes (b) to (d)
receive the ink color. Not being classified does not mean not converging:
near a pole, for example, the iterate goes very far away and takes many steps
to come back. In the 25 fractal pictures of the article, all points were
classified, except in Fig.~11: there, 174,488 of the 10,240,000
points (1.7\%) remained unclassified after 300 iterations, and the last iterate
of each of them lies within $10^{-6}$ of $0$ or $1$. Classes (b) and (d) are
empty in every picture (Table~\ref{tab:s6}). Class (b) is an artifact of the
implementation, which evaluates $p$ and $p'$ separately without cancelling
common factors; it is distinct from the poles of $N_p$, the zeros of $p'$ that
are not roots of $p$, and from the removable singularities of $p/p'$ at
multiple roots. In the magnifications of Fig.~4 (half-widths $0.300$ and $0.0375$, against $1.800$ for the full view, centered at $0$) and in Fig.~11, the grid covers the window shown, but $R$ in the
tolerance and in the classification is that of the standard window.

\emph{Roots and window.} Roots are computed with \texttt{numpy.roots}, after
removing trailing zero coefficients, and roots within $10^{-6}$ of each other
are merged. For the 38 dates, we checked in exact arithmetic that the only
possible multiple root is $z=0$, with multiplicity equal to the number of
trailing zeros of DDMMYYYY. A distinct root is called isolated if its distance
to the robust center, formed by the medians of the real and imaginary parts of
the distinct roots, exceeds three times the median distance. The window is
centered at the mean of the other distinct roots, and $R$ is 1.3 times the
largest distance from one of them to that center. Isolated roots are excluded only from the framing: they remain in the polynomial and in the iteration, and part of their basins appears in the picture. The framing rule is a visualization convention, not a property of the polynomial.

\emph{Color and brightness.} Let $k$ be the index, counted from 0, of the step
at which the criterion was met, and $\Delta$ that step. For a simple root, the
smoothed count is $\nu=k-\log_2 q$, with $q=\log|\Delta|/\log(10^{-10}R)$
clipped to $[1,2^{3/2}]$; for a root of multiplicity $m$, it is
$\nu=k-\min\bigl\{1,\log\bigl(10^{-10}R/|\Delta|\bigr)/\log\bigl(m/(m-1)\bigr)\bigr\}$.
With $\nu_0$ the second percentile of $\nu$ in the basin, the brightness
parameter is $t=1-\exp\bigl(-\max(\nu-\nu_0,0)/6\bigr)$, and the color goes
from the light tone of the basin ($t=0$) to the ink tone as $t\to1$, by
interpolation in OKLab space. Without normalization, as in
Fig.~9b, a single $\nu_0$ is used for the whole image. Basin
colors have fixed lightness and chroma in OKLCH space (0.88 and 0.075), and
their hues follow the angular order of the roots around the center.

\emph{Critical points.} The article calls \emph{free critical points} the zeros
of $p''$ that are not zeros of $p$ or of $p'$. Table~\ref{tab:s3} calls them
\emph{regular free critical points}, to distinguish them from the
\emph{critical poles}, the zeros of $p'$ of multiplicity $\mu\ge2$ that are not
roots of $p$. At a critical pole, $N_p$ has a pole of order $\mu$ and local
degree $\mu$, and it counts as $\mu-1$ critical points. A zero of multiplicity
$\mu$ of $p'$ is a zero of multiplicity $\mu-1$ of $p''$, so critical poles also
appear among the zeros of $p''$, but there the formula $N_p'=p\,p''/p'^2$ does
not hold. A multiple root of $p$ is not a critical point of $N_p$, since
$N_p'(\zeta)=1-1/m\neq0$. When the roots of $p$ are simple and $p'$ has no
multiple zeros, the count is direct: $d$ roots and $d-2$ zeros of $p''$. For
01/01/2000, $k=4$, and the six critical points are the three simple roots of
$z^3+z+2$ and the three nonzero zeros of $p''(z)=6z(5z^3+2z+2)$; the zero $z=0$
of $p''$ is the triple root, which is not critical. The only date with a
critical pole is 11/02/2016, with $p(z)=z^7+z^6+2z^4+2z^3+z+6$:
$p'(z)=(z+1)^2(z^2-z+1)(7z^2-z+1)$ and $p(-1)=5$, so that $z=-1$ is a double
pole of $N_p$, whose orbit is exactly $-1\mapsto\infty\mapsto\infty$.
The critical poles are the irreducible factors of
$\gcd(p',p'')$ that do not divide $p$; their multiplicity in $p'$ gives the
order of the pole. The regular free critical points are the zeros of the
polynomial obtained from $p''$ by removing all factors it has in common with
$p\,p'$. For each date, the total number of critical points, counted with
multiplicity, is $2k-2$. Only the regular points are iterated numerically,
starting from approximations with 30 digits, with the stopping criterion
$|\Delta|<10^{-14}\max(1,\max_j|\zeta_j|)$ and distance less than $10^{-9}$ to
a root. For 20/11/1924, $\lambda_n$ was computed in double precision
($n=10^4,10^5,10^6$), in extended precision (\texttt{long double}, $n=10^5$),
with 40 digits (\texttt{mpmath}, $n=2\times10^4$), and from $c\pm10^{-8}$ and
$c+10^{-4}$ ($n=10^5$); the runs are in Table~\ref{tab:s4}.

\emph{Why the test of the critical orbits works.} Suppose that the orbit of
every free critical point converges to a root. By the classification of Fatou
components due to Fatou and Sullivan, every Fatou component is eventually mapped
into a periodic one, and a periodic component is part of the immediate basin of
an attracting or parabolic cycle, or a rotation domain (a Siegel disk or a
Herman ring); wandering components do not exist, by Sullivan's theorem. An
attracting cycle other than the roots would contain a critical point in its
immediate basin, by Fatou's theorem (Theorem~2 of the article); that point could
be neither a root, which is fixed, nor a free critical point, whose orbit
converges to a root, nor a critical pole, which lies in the Julia set. A
parabolic cycle would likewise attract the orbit of some critical point, which
is again impossible. The boundary of a rotation domain is an infinite subset of
the Julia set contained in the closure of the critical orbits, and here that
closure meets the Julia set only in the finite set formed by the critical poles
and $\infty$. Hence the Fatou set of $N_p$ is exactly the union of the basins of
the roots. The references are those of the article: Carleson and Gamelin,
\emph{Complex Dynamics} (1993); Milnor, \emph{Dynamics in One Complex
Variable} (2006); and Sullivan, \emph{Annals of Mathematics} 122 (1985).

\emph{Isolated roots.} Of the 38 dates, 14 have $c_1/c_0\ge3.5$, and the
framing rule classifies one of their roots as isolated; 13 of them fall under
Proposition~1. The exception is 14/09/2015, whose isolated root, $-4.43$, is
real, but day 14 escapes the estimate of the proposition (Table~\ref{tab:s2}).

\emph{Mandelbrot's date.} For 20/11/1924, the real root of $p$ is $-1.3169$,
and the real poles of $N_p$ are $-0.9730$ and $-0.1130$. The intervals $K_\pm$ of Proposition~3 are not the smallest possible. Write
$x_j=N_p^j(c)$. A finer verification with interval arithmetic
(Table~\ref{tab:s5}) shows that
$I_-=[x_2,x_4]\subset K_-$ and $I_+=[x_1,x_3]\subset K_+$ satisfy
$N_p(I_-)=I_+$ and $N_p(I_+)=I_-$, and that $c\in I_-$. All iterates of $c$
therefore lie in $I_-\cup I_+$; since the endpoints of $I_-$ and $I_+$ are
themselves iterates, these are the smallest closed intervals containing the
even and the odd iterates of $c$, respectively. We call them the critical-orbit
intervals.

\emph{Critical-orbit intervals.} To show that the orbit of $c$ stays in
$I_-=[x_2,x_4]$ and $I_+=[x_1,x_3]$, we isolated $c$ by Sturm sequences in an
interval of width less than $10^{-40}$, enclosed $x_1,\dots,x_5$ in intervals of
width less than $10^{-36}$ by interval arithmetic with 80 digits
(\texttt{mpmath.iv}), converted the endpoints, without rounding, into rational
numbers, and verified by exact root counting and exact comparisons that $p$ and
$p'$ have no zeros on rational intervals $J_\pm\supset I_\pm$, that $p''$ has
no zeros on $J_+$ and only the simple zero $c$ on $J_-$, and that
$x_2<c<x_4$, $x_1<x_3$, and $x_5<x_3$. These facts give $N_p(I_+)=I_-$ by
monotonicity and $N_p(I_-)=[x_1,\max(x_3,x_5)]=I_+$, since $c$ is a minimum.
Table~\ref{tab:s5} summarizes this verification; the exact data are
in the file \texttt{results/orbit\_intervals\_certificate.json} of the code repository.

\emph{Variations.} The programs also allow variations of the construction: changing the order of the digits to YYYYMMDD;
magnifying a stretch of the real axis for a date without real roots; or
replacing the Newton step by the relaxed step $z_{k+1}=z_k-a\,p(z_k)/p'(z_k)$,
with complex $a$ close to 1, which turns the superattracting roots into
attracting fixed points with multiplier $1-a$ and distorts the basins (see the \texttt{README.md} of the code repository).



% BEGIN TABLES S1-S6 (written by make_supplement.py; do not edit by hand)
{\small
\begin{longtable}{@{}l >{\raggedright\arraybackslash}p{5.6cm} l c c c@{}}
\caption{Dates, polynomials, degree $d$, number $k$ of distinct roots, and multiplicity $m_0$ of the root $z=0$ (0 if $z=0$ is not a root). All other roots are simple (exact verification)}\label{tab:s1}\\
\toprule
Date & Occasion & $p(z)$ & $d$ & $k$ & $m_0$\\
\midrule\endfirsthead
\toprule Date & Occasion & $p(z)$ & $d$ & $k$ & $m_0$\\\midrule\endhead
25/12/1642 & Birth of Isaac Newton (Julian) & $2z^7+5z^6+z^5+2z^4+z^3+6z^2+4z+2$ & 7 & 7 & 0\\
04/01/1643 & Birth of Isaac Newton (Gregorian) & $4z^6+z^4+z^3+6z^2+4z+3$ & 6 & 6 & 0\\
15/04/1707 & Birth of Leonhard Euler & $z^7+5z^6+4z^4+z^3+7z^2+7$ & 7 & 7 & 0\\
30/04/1777 & Birth of Carl Friedrich Gauss & $3z^7+4z^4+z^3+7z^2+7z+7$ & 7 & 7 & 0\\
16/08/1821 & Birth of Arthur Cayley & $z^7+6z^6+8z^4+z^3+8z^2+2z+1$ & 7 & 7 & 0\\
29/04/1854 & Birth of Henri Poincaré & $2z^7+9z^6+4z^4+z^3+8z^2+5z+4$ & 7 & 7 & 0\\
28/02/1878 & Birth of Pierre Fatou & $2z^7+8z^6+2z^4+z^3+8z^2+7z+8$ & 7 & 7 & 0\\
03/02/1893 & Birth of Gaston Julia & $3z^6+2z^4+z^3+8z^2+9z+3$ & 6 & 6 & 0\\
20/11/1924 & Birth of Benoit Mandelbrot & $2z^7+z^5+z^4+z^3+9z^2+2z+4$ & 7 & 7 & 0\\
05/07/1687 & Publication of the Principia (Julian) & $5z^6+7z^4+z^3+6z^2+8z+7$ & 6 & 6 & 0\\
15/07/1687 & Publication of the Principia (Gregorian) & $z^7+5z^6+7z^4+z^3+6z^2+8z+7$ & 7 & 7 & 0\\
07/09/1822 & Independence of Brazil & $7z^6+9z^4+z^3+8z^2+2z+2$ & 6 & 6 & 0\\
23/09/1846 & Galle observes Neptune & $2z^7+3z^6+9z^4+z^3+8z^2+4z+6$ & 7 & 7 & 0\\
13/05/1888 & Golden Law (abolition of slavery in Brazil) & $z^7+3z^6+5z^4+z^3+8z^2+8z+8$ & 7 & 7 & 0\\
21/01/1889 & Announcement of King Oscar II's prize & $2z^7+z^6+z^4+z^3+8z^2+8z+9$ & 7 & 7 & 0\\
15/11/1889 & Proclamation of the Republic of Brazil & $z^7+5z^6+z^5+z^4+z^3+8z^2+8z+9$ & 7 & 7 & 0\\
30/11/1889 & Poincaré asks to stop the printing & $3z^7+z^5+z^4+z^3+8z^2+8z+9$ & 7 & 7 & 0\\
08/08/1900 & Hilbert's problems & $8z^6+8z^4+z^3+9z^2$ & 6 & 5 & 2\\
23/10/1906 & Flight of the 14-bis & $2z^7+3z^6+z^5+z^3+9z^2+6$ & 7 & 7 & 0\\
29/05/1919 & Sobral eclipse & $2z^7+9z^6+5z^4+z^3+9z^2+z+9$ & 7 & 7 & 0\\
15/10/1952 & Founding of IMPA & $z^7+5z^6+z^5+z^3+9z^2+5z+2$ & 7 & 7 & 0\\
29/12/1972 & Lorenz's talk & $2z^7+9z^6+z^5+2z^4+z^3+9z^2+7z+2$ & 7 & 7 & 0\\
16/05/1980 & Founding of LNCC & $z^7+6z^6+5z^4+z^3+9z^2+8z$ & 7 & 7 & 1\\
05/10/1988 & Brazilian Federal Constitution & $5z^6+z^5+z^3+9z^2+8z+8$ & 6 & 6 & 0\\
23/06/1993 & Wiles's announcement & $2z^7+3z^6+6z^4+z^3+9z^2+9z+3$ & 7 & 7 & 0\\
22/04/1500 & Arrival of Cabral (Julian) & $2z^7+2z^6+4z^4+z^3+5z^2$ & 7 & 6 & 2\\
02/05/1500 & Arrival of Cabral (proleptic Gregorian) & $2z^6+5z^4+z^3+5z^2$ & 6 & 5 & 2\\
04/10/1582 & Last Julian day & $4z^6+z^5+z^3+5z^2+8z+2$ & 6 & 6 & 0\\
15/10/1582 & First Gregorian day & $z^7+5z^6+z^5+z^3+5z^2+8z+2$ & 7 & 7 & 0\\
20/07/1969 & Apollo 11 (Brasília time) & $2z^7+7z^4+z^3+9z^2+6z+9$ & 7 & 7 & 0\\
21/07/1969 & Apollo 11 (UTC) & $2z^7+z^6+7z^4+z^3+9z^2+6z+9$ & 7 & 7 & 0\\
14/09/2015 & Detection of gravitational waves & $z^7+4z^6+9z^4+2z^3+z+5$ & 7 & 7 & 0\\
11/02/2016 & Announcement of the detection & $z^7+z^6+2z^4+2z^3+z+6$ & 7 & 7 & 0\\
01/01/2000 & New Year's Day 2000 & $z^6+z^4+2z^3$ & 6 & 4 & 3\\
30/08/1984 & Birth of the author & $3z^7+8z^4+z^3+9z^2+8z+4$ & 7 & 7 & 0\\
03/07/1986 & Birth of the author's wife & $3z^6+7z^4+z^3+9z^2+8z+6$ & 6 & 6 & 0\\
27/04/1952 & Birth of the author's father & $2z^7+7z^6+4z^4+z^3+9z^2+5z+2$ & 7 & 7 & 0\\
15/08/1950 & Birth of the author's mother & $z^7+5z^6+8z^4+z^3+9z^2+5z$ & 7 & 7 & 1\\
\bottomrule
\end{longtable}

\begin{longtable}{@{}l r r c r r r@{}}
\caption{Window (real center and half-width $R$), isolated root outside the window, ratio $c_1/c_0$, real root of largest modulus, and relative error of the prediction $-c_1/c_0$}\label{tab:s2}\\
\toprule
Date & center & $R$ & isolated & $c_1/c_0$ & real root & rel.\ error\\
\midrule\endfirsthead
\toprule Date & center & $R$ & isolated & $c_1/c_0$ & real root & rel.\ error\\\midrule\endhead
25/12/1642 & $-0.357$ & $2.772$ & no & $2.50$ & $-2.490$ & $0.004$\\
04/01/1643 & $0.000$ & $1.623$ & no & $0.00$ & --- & ---\\
15/04/1707 & $0.026$ & $1.472$ & yes & $5.00$ & $-5.154$ & $0.030$\\
30/04/1777 & $0.000$ & $1.800$ & no & $0.00$ & $-1.302$ & ---\\
16/08/1821 & $0.035$ & $1.390$ & yes & $6.00$ & $-6.209$ & $0.034$\\
29/04/1854 & $0.016$ & $1.360$ & yes & $4.50$ & $-4.597$ & $0.021$\\
28/02/1878 & $0.011$ & $1.480$ & yes & $4.00$ & $-4.065$ & $0.016$\\
03/02/1893 & $0.000$ & $1.895$ & no & $0.00$ & --- & ---\\
20/11/1924 & $0.000$ & $1.833$ & no & $0.00$ & $-1.317$ & ---\\
05/07/1687 & $0.000$ & $1.622$ & no & $0.00$ & --- & ---\\
15/07/1687 & $0.042$ & $1.602$ & yes & $5.00$ & $-5.253$ & $0.048$\\
07/09/1822 & $0.000$ & $1.324$ & no & $0.00$ & --- & ---\\
23/09/1846 & $-0.214$ & $2.810$ & no & $1.50$ & $-2.376$ & $0.369$\\
13/05/1888 & $-0.429$ & $3.918$ & no & $3.00$ & $-3.443$ & $0.129$\\
21/01/1889 & $-0.071$ & $1.960$ & no & $0.50$ & $-1.418$ & $0.647$\\
15/11/1889 & $-0.027$ & $1.676$ & yes & $5.00$ & $-4.840$ & $0.033$\\
30/11/1889 & $0.000$ & $1.813$ & no & $0.00$ & $-1.163$ & ---\\
08/08/1900 & $0.000$ & $1.377$ & no & $0.00$ & $0.000$ & ---\\
23/10/1906 & $-0.214$ & $1.970$ & no & $1.50$ & $-1.729$ & $0.133$\\
29/05/1919 & $0.020$ & $1.375$ & yes & $4.50$ & $-4.622$ & $0.026$\\
15/10/1952 & $-0.034$ & $1.596$ & yes & $5.00$ & $-4.798$ & $0.042$\\
29/12/1972 & $-0.010$ & $1.415$ & yes & $4.50$ & $-4.442$ & $0.013$\\
16/05/1980 & $0.022$ & $1.544$ & yes & $6.00$ & $-6.134$ & $0.022$\\
05/10/1988 & $-0.033$ & $1.755$ & no & $0.20$ & --- & ---\\
23/06/1993 & $-0.214$ & $2.574$ & no & $1.50$ & $-2.195$ & $0.317$\\
22/04/1500 & $-0.167$ & $2.104$ & no & $1.00$ & $-1.786$ & $0.440$\\
02/05/1500 & $0.000$ & $1.794$ & no & $0.00$ & $0.000$ & ---\\
04/10/1582 & $-0.042$ & $1.685$ & no & $0.25$ & $-0.932$ & $0.732$\\
15/10/1582 & $-0.035$ & $1.564$ & yes & $5.00$ & $-4.789$ & $0.044$\\
20/07/1969 & $0.000$ & $2.182$ & no & $0.00$ & $-1.679$ & ---\\
21/07/1969 & $-0.071$ & $2.290$ & no & $0.50$ & $-1.833$ & $0.727$\\
14/09/2015 & $0.072$ & $1.908$ & yes & $4.00$ & $-4.435$ & $0.098$\\
11/02/2016 & $-0.143$ & $1.867$ & no & $1.00$ & $-1.579$ & $0.367$\\
01/01/2000 & $0.000$ & $1.839$ & no & $0.00$ & $-1.000$ & ---\\
30/08/1984 & $0.000$ & $1.928$ & no & $0.00$ & $-1.483$ & ---\\
03/07/1986 & $0.000$ & $1.822$ & no & $0.00$ & --- & ---\\
27/04/1952 & $0.027$ & $1.422$ & yes & $3.50$ & $-3.661$ & $0.044$\\
15/08/1950 & $0.048$ & $1.542$ & yes & $5.00$ & $-5.290$ & $0.055$\\
\bottomrule
\end{longtable}

\begin{longtable}{@{}l >{\raggedright\arraybackslash}p{12.4cm}@{}}
\caption{Free critical points of $N_p$. For each regular one, $z_0$ (a zero of $p''$ that is not a zero of $p$ or of $p'$), the root reached and the number of iterations, with the criterion: step smaller than $10^{-14}\max(1,\max|\zeta|)$ and distance to a root smaller than $10^{-9}$, at most $10^4$ iterations. Critical poles (multiple zeros of $p'$ that are not roots of $p$) were identified in exact arithmetic and were not iterated numerically}\label{tab:s3}\\
\toprule
Date & $z_0\to$ root reached (iterations)\\
\midrule\endfirsthead
\toprule Date & $z_0\to$ root reached (iterations)\\\midrule\endhead
25/12/1642 & $-1.745$ $\to$ $-2.489$ (68); $-0.329-0.438i$ $\to$ $-0.401-0.484i$ (5); $-0.329+0.438i$ $\to$ $-0.401+0.484i$ (5); $0.308-0.421i$ $\to$ $-0.401-0.484i$ (7); $0.308+0.421i$ $\to$ $-0.401+0.484i$ (7)\\
04/01/1643 & $-0.366-0.387i$ $\to$ $-0.255-0.733i$ (9); $-0.366+0.387i$ $\to$ $-0.255+0.733i$ (9); $0.366-0.467i$ $\to$ $-0.255-0.733i$ (12); $0.366+0.467i$ $\to$ $-0.255+0.733i$ (12)\\
15/04/1707 & $-3.656$ $\to$ $-5.154$ (88); $-0.253-0.451i$ $\to$ $0.065-0.948i$ (10); $-0.253+0.451i$ $\to$ $0.065+0.948i$ (10); $0.296-0.503i$ $\to$ $0.065-0.948i$ (9); $0.296+0.503i$ $\to$ $0.065+0.948i$ (9)\\
30/04/1777 & $-0.793$ $\to$ $-1.302$ (12); $-0.101-0.477i$ $\to$ $-0.562-0.638i$ (8); $-0.101+0.477i$ $\to$ $-0.562+0.638i$ (8); $0.497-0.587i$ $\to$ $-0.562-0.638i$ (8); $0.497+0.587i$ $\to$ $-0.562+0.638i$ (8)\\
16/08/1821 & $-4.403$ $\to$ $-6.209$ (96); $-0.126-0.459i$ $\to$ $-0.148-0.341i$ (5); $-0.126+0.459i$ $\to$ $-0.148+0.341i$ (5); $0.184-0.591i$ $\to$ $-0.148-0.341i$ (6); $0.184+0.591i$ $\to$ $-0.148+0.341i$ (6)\\
29/04/1854 & $-3.267$ $\to$ $-4.597$ (28); $-0.265-0.395i$ $\to$ $-0.541-0.547i$ (8); $-0.265+0.395i$ $\to$ $-0.541+0.547i$ (8); $0.292-0.414i$ $\to$ $-0.541-0.547i$ (9); $0.292+0.414i$ $\to$ $-0.541+0.547i$ (9)\\
28/02/1878 & $-2.891$ $\to$ $-4.065$ (14); $-0.314-0.387i$ $\to$ $-0.738-0.576i$ (10); $-0.314+0.387i$ $\to$ $-0.738+0.576i$ (10); $0.331-0.394i$ $\to$ $-0.738-0.576i$ (8); $0.331+0.394i$ $\to$ $-0.738+0.576i$ (8)\\
03/02/1893 & $-0.381-0.485i$ $\to$ $-0.527-0.209i$ (6); $-0.381+0.485i$ $\to$ $-0.527+0.209i$ (6); $0.381-0.568i$ $\to$ $-0.527-0.209i$ (8); $0.381+0.568i$ $\to$ $-0.527+0.209i$ (8)\\
20/11/1924 & $-0.689$: no convergence in $10^4$ iterations; nonconvergence proved (Proposition~3 of the article); $-0.189-0.705i$ $\to$ $-0.100-0.687i$ (5); $-0.189+0.705i$ $\to$ $-0.100+0.687i$ (5); $0.534-0.547i$ $\to$ $-0.100-0.687i$ (6); $0.534+0.547i$ $\to$ $-0.100+0.687i$ (6)\\
05/07/1687 & $-0.100-0.436i$ $\to$ $-0.630-0.473i$ (8); $-0.100+0.436i$ $\to$ $-0.630+0.473i$ (8); $0.100-0.624i$ $\to$ $-0.630-0.473i$ (8); $0.100+0.624i$ $\to$ $-0.630+0.473i$ (8)\\
15/07/1687 & $-3.715$ $\to$ $-5.253$ (39); $-0.098-0.420i$ $\to$ $-0.625-0.479i$ (8); $-0.098+0.420i$ $\to$ $-0.625+0.479i$ (8); $0.169-0.621i$ $\to$ $-0.625-0.479i$ (8); $0.169+0.621i$ $\to$ $-0.625+0.479i$ (8)\\
07/09/1822 & $-0.114-0.455i$ $\to$ $-0.208-0.547i$ (5); $-0.114+0.455i$ $\to$ $-0.208+0.547i$ (5); $0.114-0.577i$ $\to$ $-0.208-0.547i$ (6); $0.114+0.577i$ $\to$ $-0.208+0.547i$ (6)\\
23/09/1846 & $-1.591$ $\to$ $-2.376$ (22); $-0.050-0.404i$ $\to$ $-0.500-0.630i$ (8); $-0.050+0.404i$ $\to$ $-0.500+0.630i$ (8); $0.310-0.792i$ $\to$ $0.166-1.141i$ (9); $0.310+0.792i$ $\to$ $0.166+1.141i$ (9)\\
13/05/1888 & $-2.393$ $\to$ $-3.443$ (11); $-0.181-0.529i$ $\to$ $-0.642-0.617i$ (7); $-0.181+0.529i$ $\to$ $-0.642+0.617i$ (7); $0.307-0.645i$ $\to$ $-0.642-0.617i$ (6); $0.307+0.645i$ $\to$ $-0.642+0.617i$ (6)\\
21/01/1889 & $-0.834$ $\to$ $-1.418$ (14); $-0.274-0.615i$ $\to$ $-0.650-0.812i$ (7); $-0.274+0.615i$ $\to$ $-0.650+0.812i$ (7); $0.512-0.492i$ $\to$ $-0.650-0.812i$ (9); $0.512+0.492i$ $\to$ $-0.650+0.812i$ (9)\\
15/11/1889 & $-3.457$ $\to$ $-4.840$ (21); $-0.412-0.426i$ $\to$ $-0.861-0.642i$ (10); $-0.412+0.426i$ $\to$ $-0.861+0.642i$ (10); $0.355-0.433i$ $\to$ $-0.861-0.642i$ (10); $0.355+0.433i$ $\to$ $-0.861+0.642i$ (10)\\
30/11/1889 & $-0.625$ $\to$ $-1.163$ (29); $-0.179-0.628i$ $\to$ $-0.563-0.787i$ (8); $-0.179+0.628i$ $\to$ $-0.563+0.787i$ (8); $0.492-0.484i$ $\to$ $-0.563-0.787i$ (9); $0.492+0.484i$ $\to$ $-0.563+0.787i$ (9)\\
08/08/1900 & $-0.195-0.454i$ $\to$ $0.000$ (45); $-0.195+0.454i$ $\to$ $0.000$ (45); $0.195-0.520i$ $\to$ $0.000$ (46); $0.195+0.520i$ $\to$ $0.000$ (46)\\
23/10/1906 & $-0.983$ $\to$ $-1.729$ (73); $-0.471-0.602i$ $\to$ $0.027-0.769i$ (7); $-0.471+0.602i$ $\to$ $0.027+0.769i$ (7); $0.427-0.436i$ $\to$ $0.027-0.769i$ (8); $0.427+0.436i$ $\to$ $0.027+0.769i$ (8)\\
29/05/1919 & $-3.280$ $\to$ $-4.622$ (96); $-0.256-0.417i$ $\to$ $0.023-0.901i$ (10); $-0.256+0.417i$ $\to$ $0.023+0.901i$ (10); $0.289-0.436i$ $\to$ $0.023-0.901i$ (13); $0.289+0.436i$ $\to$ $0.023+0.901i$ (13)\\
15/10/1952 & $-3.432$ $\to$ $-4.798$ (145); $-0.458-0.414i$ $\to$ $-0.283-0.405i$ (5); $-0.458+0.414i$ $\to$ $-0.283+0.405i$ (5); $0.388-0.422i$ $\to$ $-0.283-0.405i$ (7); $0.388+0.422i$ $\to$ $-0.283+0.405i$ (7)\\
29/12/1972 & $-3.167$ $\to$ $-4.442$ (147); $-0.340-0.391i$ $\to$ $-0.398-0.242i$ (5); $-0.340+0.391i$ $\to$ $-0.398+0.242i$ (5); $0.317-0.389i$ $\to$ $-0.398-0.242i$ (8); $0.317+0.389i$ $\to$ $-0.398+0.242i$ (8)\\
16/05/1980 & $-4.360$ $\to$ $-0.680$ (13); $-0.258-0.467i$ $\to$ $0.000$ (7); $-0.258+0.467i$ $\to$ $0.000$ (7); $0.295-0.508i$ $\to$ $0.000$ (7); $0.295+0.508i$ $\to$ $0.000$ (7)\\
05/10/1988 & $-0.451-0.383i$ $\to$ $-0.247-0.868i$ (26); $-0.451+0.383i$ $\to$ $-0.247+0.868i$ (26); $0.384-0.442i$ $\to$ $-0.247-0.868i$ (27); $0.384+0.442i$ $\to$ $-0.247+0.868i$ (27)\\
23/06/1993 & $-1.483$ $\to$ $-2.195$ (22); $-0.136-0.510i$ $\to$ $-0.446-0.248i$ (7); $-0.136+0.510i$ $\to$ $-0.446+0.248i$ (7); $0.342-0.634i$ $\to$ $-0.446-0.248i$ (7); $0.342+0.634i$ $\to$ $-0.446+0.248i$ (7)\\
22/04/1500 & $-1.159$ $\to$ $0.000$ (47); $-0.127-0.444i$ $\to$ $0.000$ (45); $-0.127+0.444i$ $\to$ $0.000$ (45); $0.350-0.599i$ $\to$ $0.000$ (45); $0.350+0.599i$ $\to$ $0.000$ (45)\\
02/05/1500 & $-0.080-0.443i$ $\to$ $0.000$ (45); $-0.080+0.443i$ $\to$ $0.000$ (45); $0.080-0.904i$ $\to$ $0.000$ (46); $0.080+0.904i$ $\to$ $0.000$ (46)\\
04/10/1582 & $-0.424-0.328i$ $\to$ $-0.304$ (6); $-0.424+0.328i$ $\to$ $-0.304$ (6); $0.341-0.417i$ $\to$ $-0.304$ (7); $0.341+0.417i$ $\to$ $-0.304$ (7)\\
15/10/1582 & $-3.431$ $\to$ $-0.928$ (12); $-0.402-0.336i$ $\to$ $-0.305$ (6); $-0.402+0.336i$ $\to$ $-0.305$ (6); $0.332-0.378i$ $\to$ $-0.305$ (7); $0.332+0.378i$ $\to$ $-0.305$ (7)\\
20/07/1969 & $-1.041$ $\to$ $-1.679$ (52); $-0.053-0.450i$ $\to$ $-0.537-0.704i$ (8); $-0.053+0.450i$ $\to$ $-0.537+0.704i$ (8); $0.574-0.820i$ $\to$ $0.263-1.165i$ (9); $0.574+0.820i$ $\to$ $0.263+1.165i$ (9)\\
21/07/1969 & $-1.165$ $\to$ $-1.833$ (11); $-0.064-0.463i$ $\to$ $-0.555-0.705i$ (8); $-0.064+0.463i$ $\to$ $-0.555+0.705i$ (8); $0.468-0.790i$ $\to$ $1.025-1.008i$ (11); $0.468+0.790i$ $\to$ $1.025+1.008i$ (11)\\
14/09/2015 & $-3.113$ $\to$ $-4.435$ (14); $-0.110$ $\to$ $-4.435$ (8); $0.000$ $\to$ $-4.435$ (8); $0.183-0.896i$ $\to$ $0.615+0.596i$ (10); $0.183+0.896i$ $\to$ $0.615-0.596i$ (10)\\
11/02/2016 & $-0.457$ $\to$ $-1.579$ (14); $0.000$ $\to$ $-1.579$ (17); $0.371-0.698i$ $\to$ $0.933-0.739i$ (13); $0.371+0.698i$ $\to$ $0.933+0.739i$ (13); critical pole $-1$, of order 2: exact orbit $-1\mapsto\infty\mapsto\infty$\\
01/01/2000 & $-0.560$ $\to$ $0.000$ (75); $0.280-0.797i$ $\to$ $0.000$ (77); $0.280+0.797i$ $\to$ $0.000$ (77)\\
30/08/1984 & $-0.954$ $\to$ $-1.483$ (67); $-0.049-0.421i$ $\to$ $-0.433-0.393i$ (7); $-0.049+0.421i$ $\to$ $-0.433+0.393i$ (7); $0.526-0.747i$ $\to$ $-0.433-0.393i$ (10); $0.526+0.747i$ $\to$ $-0.433+0.393i$ (10)\\
03/07/1986 & $-0.091-0.540i$ $\to$ $-0.525-0.522i$ (7); $-0.091+0.540i$ $\to$ $-0.525+0.522i$ (7); $0.091-0.811i$ $\to$ $-0.525-0.522i$ (7); $0.091+0.811i$ $\to$ $-0.525+0.522i$ (7)\\
27/04/1952 & $-2.586$ $\to$ $-3.661$ (33); $-0.277-0.441i$ $\to$ $-0.309-0.372i$ (5); $-0.277+0.441i$ $\to$ $-0.309+0.372i$ (5); $0.320-0.451i$ $\to$ $-0.309-0.372i$ (7); $0.320+0.451i$ $\to$ $-0.309+0.372i$ (7)\\
15/08/1950 & $-3.735$ $\to$ $-0.474$ (13); $-0.120-0.488i$ $\to$ $0.000$ (7); $-0.120+0.488i$ $\to$ $0.000$ (7); $0.202-0.643i$ $\to$ $0.000$ (8); $0.202+0.643i$ $\to$ $0.000$ (8)\\
\bottomrule
\end{longtable}

\begin{table}[htbp]
\centering
\caption{Orbit of the real free critical point $c\approx-0.688805$ of $N_p$ for 20/11/1924: indicator $\lambda_n=\frac1n\sum_{j=1}^n\log|N_p'(x_j)|$,
with $x_j=N_p^j(c)$, and interval visited by the orbit (minimum and maximum of the iterates)}
\label{tab:s4}
\begin{tabular}{@{}l l r r l@{}}
\toprule
Precision & starting point & $n$ & $\lambda_n$ & visited interval\\
\midrule
double & $c$ & $10^{4}$ & $0.1917$ & [$-0.7622$, $0.2005$]\\
double & $c$ & $10^{5}$ & $0.1830$ & [$-0.7622$, $0.2005$]\\
double & $c$ & $10^{6}$ & $0.1854$ & [$-0.7622$, $0.2005$]\\
extended (\texttt{long double}) & $c$ & $10^{5}$ & $0.1849$ & [$-0.7622$, $0.2005$]\\
double & $c+10^{-8}$ & $10^{5}$ & $0.1835$ & [$-0.7622$, $0.2005$]\\
double & $c-10^{-8}$ & $10^{5}$ & $0.1835$ & [$-0.7622$, $0.2005$]\\
double & $c+10^{-4}$ & $10^{5}$ & $0.1845$ & [$-0.7622$, $0.2005$]\\
40 digits (\texttt{mpmath}) & $c$ & $2\times10^{4}$ & $0.1835$ & [$-0.7622$, $0.2005$]\\
\bottomrule
\end{tabular}
\end{table}

\begin{table}[htbp]
\centering
\caption{Critical-orbit intervals for 20/11/1924: $I_-=[x_2,x_4]$ and $I_+=[x_1,x_3]$, the
smallest closed intervals that contain the even and the odd iterates of $c$, respectively.
Validated enclosures of $x_j=N_p^j(c)$ ($x_0=c$), given by midpoint and radius, and result of
each exact check. The exact rational endpoints, the isolating intervals, the Sturm counts, and
the bounds for each inequality are in the file
\texttt{results/orbit\_intervals\_certificate.json} of the code repository}
\label{tab:s5}
\begin{tabular}{@{}c@{}}
\begin{tabular}{@{}l r r@{}}
\toprule
 & midpoint & radius\\
\midrule
$x_0$ & $-0.688805042592084$ & $2.2\times10^{-41}$\\
$x_1$ & $0.158090063368374$ & $1.2\times10^{-40}$\\
$x_2$ & $-0.762158290647476$ & $6.9\times10^{-40}$\\
$x_3$ & $0.200462120038431$ & $5.0\times10^{-39}$\\
$x_4$ & $-0.626668751444423$ & $2.4\times10^{-38}$\\
$x_5$ & $0.179068904269946$ & $1.2\times10^{-37}$\\
\bottomrule
\end{tabular}\\[8pt]
\begin{tabular}{@{}l l@{}}
\toprule
Check & result\\
\midrule
$\gcd(p',p'')=1$: $N_p$ has no critical poles & yes\\
$p''$ squarefree, a single real zero, $c$ isolated & yes\\
on $J_-\supset I_-$: $p$ and $p'$ without zeros, $p''$ with one zero ($c$) & yes\\
on $J_+\supset I_+$: $p$, $p'$, and $p''$ without zeros & yes\\
$x_2<c<x_4$, $x_1<x_3$, $x_2<x_4$ & yes\\
$x_5<x_3$ & yes\\
\bottomrule
\end{tabular}
\end{tabular}
\end{table}

\begin{longtable}{@{}l r r r r@{}}
\caption{Classification of the $3200\times3200$ points of each fractal picture of the article:
(a) classified; (b) non-finite value in floating point; (c) unclassified after 300 iterations;
(d) stopped by the step criterion farther than $10^{-4}R$ from every root. For Fig.~11, all 174,488 points of class (c) have their last iterate within $10^{-6}$ of $0$ or $1$}
\label{tab:s6}\\
\toprule
Picture & (a) & (b) & (c) & (d)\\
\midrule\endfirsthead
\toprule Picture & (a) & (b) & (c) & (d)\\\midrule\endhead
25/12/1642 (Fig. 1a) & 10,240,000 & 0 & 0 & 0\\
04/01/1643 (Fig. 1b) & 10,240,000 & 0 & 0 & 0\\
30/04/1777 (Figs. 4a, 5) & 10,240,000 & 0 & 0 & 0\\
15/04/1707 (Fig. 7a) & 10,240,000 & 0 & 0 & 0\\
16/08/1821 (Fig. 7b) & 10,240,000 & 0 & 0 & 0\\
29/04/1854 (Fig. 7c) & 10,240,000 & 0 & 0 & 0\\
28/02/1878 (Fig. 7d) & 10,240,000 & 0 & 0 & 0\\
03/02/1893 (Fig. 7e) & 10,240,000 & 0 & 0 & 0\\
20/11/1924 (Fig. 7f) & 10,240,000 & 0 & 0 & 0\\
05/07/1687 (Fig. 8a) & 10,240,000 & 0 & 0 & 0\\
23/09/1846 (Fig. 8b) & 10,240,000 & 0 & 0 & 0\\
30/11/1889 (Fig. 8c) & 10,240,000 & 0 & 0 & 0\\
29/05/1919 (Fig. 8d) & 10,240,000 & 0 & 0 & 0\\
29/12/1972 (Fig. 8e) & 10,240,000 & 0 & 0 & 0\\
16/05/1980 (Fig. 8f) & 10,240,000 & 0 & 0 & 0\\
08/08/1900 (Fig. 9) & 10,240,000 & 0 & 0 & 0\\
30/08/1984 (Fig. 10a) & 10,240,000 & 0 & 0 & 0\\
03/07/1986 (Fig. 10b) & 10,240,000 & 0 & 0 & 0\\
27/04/1952 (Fig. 10c) & 10,240,000 & 0 & 0 & 0\\
15/08/1950 (Fig. 10d) & 10,240,000 & 0 & 0 & 0\\
$z^2-1$ (Fig. 3a) & 10,240,000 & 0 & 0 & 0\\
$z^3-1$ (Fig. 3b) & 10,240,000 & 0 & 0 & 0\\
$z^3-2z+2$ (Fig. 11) & 10,065,512 & 0 & 174,488 & 0\\
Gauss, magnification (Fig. 4b) & 10,240,000 & 0 & 0 & 0\\
Gauss, magnification (Fig. 4c) & 10,240,000 & 0 & 0 & 0\\
\bottomrule
\end{longtable}
}
% END TABLES S1-S6

\end{document}
